\chapter{Repository ed entrypoint}
\label{chap:repository}

\section{Punti di ingresso}

Per l'uso operativo, \filepath{infer.py} riceve direttamente la cartella delle
immagini, ricostruisce tutti i nodi trovati ed esporta il risultato per QGIS:

\begin{lstlisting}[language={}]
python infer.py D:/percorso/immagini
\end{lstlisting}

Non richiede label o annotazioni. L'export usa la geometria VGGT quando
disponibile e ripiega automaticamente sul layout CAD. \filepath{run.py} rimane
l'entrypoint tecnico
per avviare la sola ricostruzione usando la configurazione del progetto:

\begin{lstlisting}[language={}]
python run.py
\end{lstlisting}

Senza argomenti vengono elaborati tutti i nodi ammessi dalla configurazione;
gli eventuali ID servono soltanto a limitare la lavorazione, come descritto nel
Capitolo~\ref{chap:esecuzione}. \filepath{run.py} mantiene volutamente un ruolo
ridotto: interpreta la richiesta e delega la ricostruzione a
\code{run_multiview_all()} o \code{run_multiview_pipeline()}.

Valutazione, visualizzazione aggiuntiva ed export QGIS sono operazioni separate
e non sono necessarie per ottenere il grafo ricostruito.

\section{Struttura del repository}

La tabella riporta i percorsi utili per orientarsi nel codice, senza elencare
ogni file di servizio presente nel repository.

\begin{contracttable}
\filepath{infer.py} & Entrypoint operativo & Riceve le immagini, avvia la
ricostruzione ed esporta il GeoJSON per QGIS. \\

\filepath{run.py} & Entrypoint tecnico & Avvia la sola ricostruzione completa
o quella dei soli ID indicati. \\

\filepath{graph_reconstruction/config.py} & Configurazione & Carica
\filepath{.env} e definisce percorsi, soglie e vincoli globali. \\

\filepath{graph_reconstruction/pipeline/} & Pipeline & Contiene
l'orchestrazione e le fasi di detection, ReID, supporto VGGT e costruzione del
grafo. \\

\filepath{graph_reconstruction/shared/},
\filepath{graph_reconstruction/visualization/} & Moduli condivisi & Raccolgono
funzioni di supporto, accesso ai dati e visualizzazioni usate dalla pipeline. \\
\midrule

\filepath{models/}, \filepath{graph_reconstruction/assets/},
\filepath{data/} & Risorse & Ospitano checkpoint, primitive grafiche e la
struttura di immagini, label e annotazioni; i percorsi effettivi sono stabiliti
dalla configurazione. \\

\filepath{output/} & Artefatti & Destinazione predefinita di risultati, cache
e file diagnostici. \\
\end{contracttable}

\subsection{Architettura funzionale}

La pipeline è prevalentemente funzionale: l'orchestratore coordina funzioni
che trasformano dataclass, dizionari e grafi NetworkX. Il flusso di chiamata è
descritto nel Capitolo~\ref{chap:flusso-codice}, mentre gli oggetti scambiati tra le
fasi sono raccolti nel Capitolo~\ref{chap:oggetti}.

\section{Strumenti successivi alla ricostruzione}

Gli script seguenti leggono gli artefatti già prodotti dalla pipeline. Sono
utili per analisi o consegne specifiche, ma non fanno parte dell'esecuzione
necessaria a ottenere il grafo.

\subsection{Valutazione}

\filepath{evaluate.py} usa \filepath{scene_nodes.json},
\filepath{fused_graph.json}, la cache delle detection e, quando disponibili,
le label estese e le annotazioni di riferimento. Calcola le metriche sulle
istanze e sugli archi e avvia il confronto con gli standard. L'assenza dei dati
di riferimento limita la valutazione, non la ricostruzione.

\subsection{Visualizzazione ed export QGIS}

\filepath{visualize.py} legge \filepath{fused_graph.json} e produce
\filepath{fused_graph_plot.png}; può inoltre generare un GeoJSON sintetico e il
report del best match. \filepath{graph_reconstruction/qgis.py} gestisce invece
l'export QGIS dedicato, per un singolo grafo o per tutti quelli disponibili,
usando geometria VGGT quando richiesta e disponibile oppure coordinate
sintetiche. Nessuno dei due modifica il JSON ricostruito.

\subsection{Diagnostica e manutenzione}

Le utility in \filepath{tools/debug/} servono per ispezioni mirate. Gli script
in \filepath{tools/benchmark/} controllano invece requisiti e consumo di
risorse. Non sono dipendenze operative del flusso normale.

\section{Caricamento della configurazione}

\filepath{graph_reconstruction/config.py} viene importato prima dell'avvio
della pipeline. In questa fase legge il file \filepath{.env} nella root senza
sovrascrivere eventuali valori già presenti nell'ambiente di processo: una
variabile impostata dall'esterno ha quindi la precedenza.

La funzione \code{_path_setting()} mantiene i percorsi assoluti e risolve quelli
relativi rispetto alla root del progetto. Se manca un percorso obbligatorio,
solleva un \code{RuntimeError} durante l'import, prima di \code{main()}. La mappa
delle impostazioni operative e dei parametri globali è riportata nel
Capitolo~\ref{chap:esecuzione} e non viene duplicata qui.
